% Dinic Algorithm

Max flow algorithm running in $O(V^2 E)$ in the worst case, but significantly faster in practice, particularly on unit networks ($O(E \sqrt{V})$).
\begin{lstlisting}
const long long INF = 1e18;

struct Dinic {
    struct Edge { int to; long long flow, cap; int rev; };
    int n, source, sink;
    vector<int> level, ptr;
    vector<vector<Edge>> adj;

    Dinic(int n, int source, int sink) : n(n), source(source), sink(sink), level(n + 1), ptr(n + 1), adj(n + 1) {}

    void add_edge(int from, int to, long long cap) {
        adj[from].push_back({to, 0, cap, (int)adj[to].size()});
        adj[to].push_back({from, 0, 0, (int)adj[from].size() - 1});
    }

    bool bfs() {
        fill(level.begin(), level.end(), -1);
        level[source] = 0;
        queue<int> q; q.push(source);
        while (!q.empty()) {
            int u = q.front(); q.pop();
            for (auto& edge : adj[u]) {
                if (edge.cap - edge.flow > 0 && level[edge.to] == -1) {
                    level[edge.to] = level[u] + 1;
                    q.push(edge.to);
                }
            }
        }
        return level[sink] != -1;
    }

    long long dfs(int u, long long pushed) {
        if (pushed == 0 || u == sink) return pushed;
        for (int& cid = ptr[u]; cid < adj[u].size(); ++cid) {
            auto& edge = adj[u][cid];
            int tr = edge.to;
            if (level[u] + 1 != level[tr] || edge.cap - edge.flow == 0) continue;
            long long tr_pushed = dfs(tr, min(pushed, edge.cap - edge.flow));
            if (tr_pushed == 0) continue;
            edge.flow += tr_pushed;
            adj[tr][edge.rev].flow -= tr_pushed;
            return tr_pushed;
        }
        return 0;
    }

    long long max_flow() {
        long long flow = 0;
        while (bfs()) {
            fill(ptr.begin(), ptr.end(), 0);
            while (long long pushed = dfs(source, INF)) {
                flow += pushed;
            }
        }
        return flow;
    }
};
\end{lstlisting}
